\section{Short-Term Projects (Summer/Winter)}

        \begin{onecolentry}
            \textbf{$\bullet$ Mathematical Modelling of Robert’s Linkage for Seismic Isolation in LIGO}, May - Jul 23.  \href{https://drive.google.com/file/d/1-pubIKgLbrkpFbzaoyhea7KZnE0AYM1R/view?usp=sharing}{\faExternalLink} \\
            \textit{Guide : Dr. Suresh Doravari, Inter-University Centre of Astronomy and Astrophysics (IUCAA), India.}
        \end{onecolentry}
        
        \vspace{0.1 cm}
        
        \begin{onecolentry}
            \textbf{$\bullet$ Study of Advanced Microscopy Techniques} (FCS, AFM), May - Jul 22.\\
            \textit{Guide : Dr. Shivprasad Patil, Indian Institute of Science Education and Research (IISER-Pune), India.}
        \end{onecolentry}    
        
        \vspace{0.1 cm}
        
        \begin{onecolentry}
            \textbf{$\bullet$ Introduction to Elementary Particles and Standard Model of Particle Physics}, Feb - Apr 2021.
            \href{https://drive.google.com/file/d/1G1gYijihDofVz5WPNy9w-gBTX9VIDLJ8/view?usp=sharing}{(B.Sc. Project )}\\
            \textit{Guide : Dr. L. M. Chaudhary, Nowrosjee Wadia College, Savitribai Phule Pune University (SPPU), India.}            
        \end{onecolentry}
        
        \vspace{0.1 cm}
        
        \begin{onecolentry}
           \textbf{$\bullet$ Innovative Laser Microscope Design}, Feb 20.\\
            \textit{Guide : Dr. Vivek Antad, Nowrosjee Wadia College (NWC), Savitribai Phule Pune University (SPPU), India.}
        \end{onecolentry}
        
        % \begin{twocolentry}{
        %     May - July 23
        % }
        %     \textbf{Mathematical Modelling of Robert’s Linkage for Seismic Isolation in LIGO}
        %     \href{https://drive.google.com/file/d/1-pubIKgLbrkpFbzaoyhea7KZnE0AYM1R/view?usp=sharing}{\faExternalLink}\end{twocolentry}

        % \vspace{0.10 cm}
        % \begin{onecolentry}
        %     \textit{Guide : Dr. Suresh Doravari, Inter-University Centre of Astronomy and Astrophysics (IUCAA), India}

            % \vspace{0.10 cm}
            % \begin{highlights}
            %     % \item \textbf{Skills Developed :} Python, Mathemtica, Linkage software.
            %     \item Modeled Roberts Linkage (RL) seismic isolator by deriving the \textit{\textbf{coupler curve equation}} and solving them numerically using \textit{\textbf{Mathematica}} and \textit{\textbf{Python}}. Used \textbf{\textit{Linkage}} software to simulate the results. Successfully reduced motion in RL due to seismic activity, contributed to the fabrication of RL for \textit{\textbf{LIGO India - IUCAA Instrumentation training cell}}.
                % Analyzed equations numerically in \textit{\textbf{Python}} and \textit{\textbf{Linkage}} Software simulations.
                % \item Explored Roberts Linkage’s application in isolating \textbf{\textit{LIGO}} detectors from seismic noise.
                % \vspace{0.10 cm}
                
            % \end{highlights}
           
        % \end{onecolentry}

   
        % \vspace{0.2 cm}

        % \begin{twocolentry}{
        %     May - July 22
        % }
        %     \textbf{Study of Advanced Microscopy Techniques} (FCS, AFM)
        %     \end{twocolentry}

        % \vspace{0.10 cm}
        % \begin{onecolentry}
        % \textit{Guide : Dr. Shivprasad Patil, Indian Institute of Science Education and Research (IISER-Pune), India}
        % \vspace{0.10 cm}
        % \begin{highlights}
        %     % \item \textbf{Tools Used :} He-Ne Laser, Avalanche Photodiode (APD), Lenses etc. \\

        %     \item Studied advanced microscopy techniques used in soft matter research including \textit{\textbf{Fluorescence Correlation Spectroscopy (FCS)}} and \textit{\textbf{Atomic Force Microscopy (AFM)}}. Assembled FCS in the lab, gained hands-on experience with FCS and AFM.
        % \end{highlights}
        
        
        % \end{onecolentry}

        % \vspace{0.2 cm}

        % \begin{twocolentry}{
        %     Feb – Apr 2021
        % }
        %     \textbf{Introduction to Elementary Particles and Standard Model of Particle Physics}
        %     \href{https://drive.google.com/file/d/1G1gYijihDofVz5WPNy9w-gBTX9VIDLJ8/view?usp=sharing}{(B.Sc. Project )}\end{twocolentry}

        % \vspace{0.10 cm}
        % \begin{onecolentry}
        % \textit{Guide : Dr. L. M. Chaudhary, Nowrosjee Wadia College, Savitribai Phule Pune University (SPPU), India}
        
        % \vspace{0.10 cm}

        % \begin{highlights} 
        %     % \item \textbf{Tools Used :} Microsoft Office
        %     \item Developed educational material to simplify concepts in particle physics and the development in the Standard Model.
        %     %Description : Developed educational materials to simplify particle physics concepts and the Standard Model during the COVID-19 lockdown, inspired by the Muon g-2 Experiment, gaining foundational knowledge in high-energy and particle physics and fostering an interest in high-energy astrophysics and astroparticle physics.
        % \end{highlights} 
        % \end{onecolentry}

        % \vspace{0.2 cm}
        
        % \begin{twocolentry}{
        %     Feb 20
        % }
        %     \textbf{Innovative Laser Microscope Design} \end{twocolentry}

        % \vspace{0.10 cm}
        % \begin{onecolentry}
        % \textit{Guide : Dr. Vivek Antad, Nowrosjee Wadia College (NWC), Savitribai Phule Pune University (SPPU), India}

        % \vspace{0.10 cm}

        % \begin{highlights}
        %     % \item \textbf{Tools Used :} Laser, Lenses, Thin film with specimen, Screen etc. 
        %     \item Designed and built a \textit{\textbf{laser microscope}} to study the optical properties of liquid droplets. Tested \textit{\textbf{optical principles}} like the lens maker’s equation, Snell’s law, Rayleigh’s criterion, magnification, refraction, and beam divergence and gained first hand experience in experimental research.
        % \end{highlights}
              
        % \end{onecolentry}